\SectionFrame{2. Základné pojmy}{Pojmy, ktoré budeme používať pri práci s datasheetom a registrami.}

\begin{frame}{Mikroprocesor verzus mikrokontrolér}
\small
\begin{tabular}{p{0.27\textwidth}p{0.31\textwidth}p{0.33\textwidth}}
\toprule
Vlastnosť & Mikroprocesor & Mikrokontrolér \\
\midrule
Hlavná úloha & všeobecné výpočty & riadenie konkrétneho zariadenia \\
Pamäť & často externá & typicky integrovaná \\
Periférie & externé obvody & GPIO, časovače, ADC, komunikácie \\
Operačný systém & často áno & často bez OS \\
Typické kritériá & výkon, pamäť, OS & cena, spotreba, real-time, spoľahlivosť \\
Príklad & PC procesor & ATmega328P, STM32, PIC, ESP32 \\
\bottomrule
\end{tabular}
\end{frame}

\begin{frame}{Čo je vo vnútri mikrokontroléra}
\begin{columns}[T,onlytextwidth]
\column{0.48\textwidth}
\begin{itemize}
  \item \textbf{CPU}: vykonáva inštrukcie.
  \item \textbf{Flash}: programová pamäť.
  \item \textbf{SRAM}: premenné a zásobník.
  \item \textbf{EEPROM}: trvalo uložené konfiguračné údaje.
  \item \textbf{Periférie}: porty, časovače, ADC, komunikácie.
\end{itemize}
\column{0.47\textwidth}
\includegraphics[width=\linewidth,height=5.0cm,keepaspectratio]{microcontroller_components.png}
\end{columns}
\end{frame}

\begin{frame}{Základné pojmy procesora}
\begin{columns}[T,onlytextwidth]
\column{0.49\textwidth}
\begin{itemize}
  \item \textbf{Inštrukcia}: elementárny príkaz CPU.
  \item \textbf{Inštrukčný súbor}: dostupné príkazy procesora.
  \item \textbf{ALU}: aritmeticko-logická jednotka.
  \item \textbf{Register}: rýchla pamäťová bunka v CPU alebo periférii.
\end{itemize}
\column{0.47\textwidth}
\begin{itemize}
  \item \textbf{PC}: programový čítač.
  \item \textbf{SP}: ukazovateľ zásobníka.
  \item \textbf{SREG}: stavový register s príznakmi.
  \item \textbf{Zásobník}: dočasné dáta, návratové adresy, obsluha prerušení.
\end{itemize}
\end{columns}
\end{frame}

\begin{frame}{Zbernice: dáta, adresy a riadenie}
\begin{center}
\begin{tikzpicture}[x=1cm,y=1cm]
\node[draw=MainBlue,fill=MainBlue,text=white,rounded corners=2pt,minimum width=2.8cm,minimum height=0.9cm] (cpu) at (0,1.8) {CPU};
\node[draw=MainBlue,fill=LightBlue,rounded corners=2pt,minimum width=2.8cm,minimum height=0.9cm] (mem) at (5.2,2.7) {pamäť};
\node[draw=MainBlue,fill=LightBlue,rounded corners=2pt,minimum width=2.8cm,minimum height=0.9cm] (io) at (5.2,0.9) {periféria};
\draw[<->,MainBlue,line width=1.1pt] (cpu) -- node[above,sloped]{dátová zbernica} (mem);
\draw[->,MainBlue,line width=1.1pt] (cpu) -- node[below,sloped]{adresná zbernica} (mem);
\draw[<->,MainBlue,line width=1.1pt] (cpu) -- node[below,sloped]{riadiace signály} (io);
\end{tikzpicture}
\end{center}
\begin{itemize}
  \item Dátová zbernica prenáša hodnoty.
  \item Adresná zbernica určuje, s ktorou bunkou alebo registrom sa pracuje.
  \item Riadiace signály určujú čítanie, zápis, prerušenie alebo stav periférie.
\end{itemize}
\end{frame}

\begin{frame}{Von Neumannova a Harvardská architektúra}
\begin{columns}[T,onlytextwidth]
\column{0.48\textwidth}
\BlueBox{Von Neumann}{
\begin{itemize}
  \item spoločná pamäť pre program aj dáta,
  \item jednoduchší model,
  \item možné úzke hrdlo pri prístupe do pamäte.
\end{itemize}}
\column{0.48\textwidth}
\BlueBox{Harvard}{
\begin{itemize}
  \item oddelená pamäť programu a dát,
  \item paralelnejší prístup,
  \item typické pre mnohé mikrokontroléry vrátane AVR.
\end{itemize}}
\end{columns}
\end{frame}

\begin{frame}{CISC, RISC a praktický význam pre C program}
\begin{itemize}
  \item \textbf{CISC}: rozsiahly súbor inštrukcií, často zložitejšie inštrukcie.
  \item \textbf{RISC}: jednoduchšie inštrukcie, viac registrov, dôraz na rýchle vykonávanie.
  \item AVR je 8-bitová RISC architektúra s 32 všeobecnými pracovnými registrami.
  \item Aj keď budeme písať v C, prekladač generuje konkrétne inštrukcie a prístupy do registrov.
\end{itemize}
\begin{alertblock}{Dôležité}
C kód v embedded systéme nie je izolovaný od hardvéru. Každý prístup do registra môže zmeniť správanie periférie.
\end{alertblock}
\end{frame}

\begin{frame}{Hodinový cyklus, strojový cyklus, real-time}
\begin{columns}[T,onlytextwidth]
\column{0.55\textwidth}
\begin{itemize}
  \item Hodinový signál určuje základný rytmus procesora.
  \item ATmega328P na Arduino UNO R3 je typicky taktovaný na 16 MHz.
  \item Nie každá inštrukcia znamená jeden užitočný regulačný výpočet.
  \item V embedded systéme musíme rozumieť časovaniu: oneskorenie, latencia prerušenia, perióda vzorkovania.
\end{itemize}
\column{0.40\textwidth}
\includegraphics[width=\linewidth,height=4.2cm,keepaspectratio]{pipeline.png}
\end{columns}
\end{frame}
